% Exercice à caractère documentaire
% Le sel de la vie
\textbf{Lire p. 34 répondre aux questions
}
\begin{enumerate}
    \item Pourquoi le sel fut-il durant plusieurs siècles une denrée très
          convoitée~?
    \item Quels facteurs permettent le dépôt du sel dans les salines~? Sur
          quelle propriété physique repose le procédé utilisé~?
    \item La sylvinite est un minerai constitué de chlorure de sodium, mélangé
          au chlorure de potassium~; elle est extraite en Alsace. La séparation
          des deux espèces dont les usages sont différents est une nécessité.
          
          Le graphique ci-dessous représente les variations de la solubilité en
          fonction de la température du chlorure de potassium et du chlorure de
          sodium~; la solubilité est égale ici à la masse maximale (en
          \unit{\g}) de soluté que l'on peut dissoudre dans $\qty{100}{\g}$
          d'eau.

          \begin{minipage}{\linewidth}
              \begin{figure}[H]
              \centering
              \includegraphics[width=\textwidth]{2025-12-09_013335-}
          \end{figure}
          \end{minipage}
            
          Dans l'industrie, de l'eau est mise en contact à
          $\qty{100}{\degreeCelsius}$ avec le minerai brut. Elle se sature en
          chlorure de sodium et en chlorure de potassium. Cette eau est ensuite
          ramenée à $\qty{10}{\degreeCelsius}$. Montrer que l'on parvient ainsi
          à séparer les deux chlorures. On déterminera à partir des courbes les
          compositions des solutions à $\qty{100}{\degreeCelsius}$ d'une part et
          à $\qty{10}{\degreeCelsius}$ d'autre part.
    
          Rechercher les usages du chlorure de potassium.
\end{enumerate}

